\section{Threat identification}
\label{chapter5}

\subsection{Identified threats and countermeasures}
\begin{enumerate}
  \item \hl{Threat 1: Eavesdropping on the Bluetooth connection.}\\
  \hl{An attacker within radio range records the Bluetooth traffic between the ESP32
  and the Raspberry Pi and reads intake times, target weights and intake results.
  These are sensitive health data. To prevent this, the two devices perform an
  ECDHE key exchange each time a connection is established and encrypt all
  messages with AES-GCM using the resulting session key (see module
  ``Bluetooth communication and encryption'' in Section 2.1). As a new session
  key is agreed for every connection, the recorded traffic does not reveal
  the content of the messages.}\\
  \hl{Requirement see} \ref{req.4.1}, \ref{req.4.2}

  \item \hl{Threat 2: Manipulation or replay of Bluetooth messages.}\\
  \hl{An attacker modifies, injects or replays messages. A modified schedule or
  target weight would lead to a wrong dose evaluation. A replayed result or
  heartbeat could hide a missed dose or an ESP32 without power and would
  suppress the Telegram warning (see Hazard 3). Since all messages are
  encrypted with AES-GCM, modified messages are detected by the authentication
  tag and are rejected. In addition, every message carries a sequence number, so
  that replayed messages are recognised and discarded (see module
  ``Bluetooth communication and encryption'' in Section 2.1).}\\
  \hl{Requirement see} \ref{req.4.2}, \ref{req.4.3}

  \item \hl{Threat 3: Disclosure of health data via the log output.}\\
  \hl{The log output of the ESP32 and the Raspberry Pi can be read by anyone with
  access to the serial interface of the ESP32 or the terminal of the Raspberry Pi. If it contained intake times,
  doses or weights, these health data would be disclosed. Therefore, the module
  ``System start and logging'' only outputs status messages (start-up,
  connection state, faults) without health data.}\\
  \hl{Requirement see} \ref{req.4.4}

  \item \hl{Threat 4: Disruption of the Bluetooth connection (jamming,
  connection drops).}\\
  \hl{An attacker deliberately interferes with the radio channel or provokes
  repeated connection drops, so that the ESP32 and the Raspberry Pi can no
  longer communicate. This corresponds to the loss of connection described in
  Hazard 3 and is handled by the same countermeasures: if the connection is lost
  during the day, the ESP32 switches to autonomous mode. There it continues to
  evaluate the intake with the schedule received at start-up, displays the
  result via the RGB LED, buffers the results and keeps trying to reconnect in
  the background (see Figure 3). The Raspberry Pi detects the disruption by the
  missing heartbeat (every 30 s) and sends a one-time Telegram alert after 2 min
  (see Figure 4). The safety function, i.e. the feedback via the LED, is
  therefore maintained even when the connection is disrupted.}\\
  \hl{Requirement see} \ref{req.3.1}--\ref{req.3.5}

  \item \hl{Threat 5: Manipulation of the scale outside the intake window.}\\
  \hl{A person or an attacker with access to the container places objects on the
  load cell or removes content in order to falsify the evaluation, or damages
  the load cell or its wiring. Such interference outside the intake window is
  detected. Outside an intake window, the measured weight is compared with the
  expected weight from the table of expected weights. Any deviation beyond the
  tolerance is detected as an unscheduled removal: the LED shows red and the
  Raspberry Pi sends a warning via Telegram. Larger changes are detected as part
  of Hazard 2: a weight change of more than 5 g outside an
  intake window as well as implausible values (raw value at the limit of the
  24-bit range, weight below {$-5$} g or above 100 g) are
  considered a faulty measurement. After two further faulty measurements, the
  ESP32 suspends the dose evaluation, the red LED flashes and the Raspberry Pi
  sends an alert via Telegram. The sensor is checked again every 10 s.}\\
  \hl{Requirement see} \ref{req.2.4}--\ref{req.2.7}, \ref{req.2.13}
  \item \hl{Threat 6: Connection to a foreign device (impersonation, man-in-the-middle).}\\
  \hl{The ECDHE key exchange alone protects against passive eavesdropping, but not
  against an attacker who actively intervenes while a connection is being
  established. A foreign device could impersonate the Raspberry Pi and transmit a
  modified schedule, or agree on a separate key with each side and thereby read
  and modify all messages (see Threats 1 and 2). To prevent this, both devices
  only communicate with their known partner. The Raspberry Pi only connects to the
  BLE address of the ESP32 stored in its configuration file, and the ESP32 only
  accepts one connection at a time. In addition, both devices hold a common
  pre-shared key (PSK), which is stored once during commissioning. The session key
  is derived from the ECDHE shared secret and the PSK. A device that does not know
  the PSK therefore derives a different session key, its messages fail the
  AES-GCM authentication and the connection is terminated. Since the BLE address
  alone can be spoofed, the PSK is the decisive protection. The ESP32 does not
  use TLS or X.509 certificates. As the PSK has no period of validity, the ESP32
  can establish the secure connection without a valid time. The current time is
  only received via this authenticated connection.}\\
  \hl{Requirement see} \ref{req.4.7}--\ref{req.4.9}

  \item \hl{Threat 7: Connection to a forged Telegram server.}\\
  \hl{An attacker in the network redirects the HTTPS connection of the Raspberry
  Pi to a forged server in order to read the reminders and warnings or to
  suppress them. Therefore, the Raspberry Pi verifies the server certificate of
  the Telegram Bot API for every HTTPS connection. This check requires a valid
  time. Before the first NTP synchronisation after start-up (e.g.\ for the alert
  of Hazard 1), the Raspberry Pi uses the last system time saved at shutdown,
  which is sufficient to check the period of validity of the certificate.}\\
  \hl{Requirement see} \ref{req.4.10}
\end{enumerate}

\subsection{Identified threats without countermeasures}
\begin{enumerate}
  \item \hl{Threat 1: Physical access to the devices.}\\
  \hl{A person with physical access to the devices can read the schedule on
  the ESP32, read the pre-shared key, read or modify the configuration file with the
  intake schedule or read the Telegram bot token on the Raspberry Pi, exchange the tablets
  in the container or tamper with the load cell. Protection of the stored data
  (e.g. flash encryption, secure boot) and a tamper-proof housing are not part of
  this project. The system is intended for use in the private home of the
  person, so physical access is assumed to be controlled by the person.}

  \item \hl{Threat 2: Exposure of health data on the Telegram side.}\\
  \hl{The reminders contain the names and quantities of the medications to be
  taken, as the person needs this information to take the correct tablets from
  the shared container. The connection to the Telegram Bot API is
  protected by HTTPS, but the messages are not end-to-end encrypted and are
  processed by an external service. Furthermore, a third party who gains
  access to the mobile phone or the Telegram account of the person can read the
  messages. The system has no influence on this and cannot detect it. The
  messages are therefore limited to the information necessary for the
  reminder.}

  \item \hl{Threat 3: Manipulated time through forged NTP responses.}\\
  \hl{NTP is an unauthenticated protocol. An attacker in the same network can
  provide the Raspberry Pi with a wrong but formally valid time. The reminders
  would then be sent at the wrong time, and the ESP32 would determine the
  intake window based on the wrong time, as the Raspberry Pi transmits the
  current time to the ESP32 each time a connection is established (see
  Requirement \ref{req.1.4} and Figure 4). The countermeasures of Hazard 1 only take effect
  if no time is received, not if the time is forged. Protection, e.g. via NTS
  or several independent time sources, is not part of this project.}

  \item \hl{Threat 4: Circumvention of the verification within the intake window.}\\
  \hl{The evaluation by {\texttt{dose\_evaluate()}} compares the weight change with
  the target weight. Within the intake window, the system therefore cannot
  distinguish whether the tablets were actually taken, were merely removed, or
  whether the weight change was caused by other objects. The scale verifies the
  removal, not the intake (see Section 1.2). Additional detection, e.g. through
  further sensors, is not part of this project.}
\end{enumerate}